% GENERATED FILE. Edit canonical lessons, then run npm run generate:books.
% canonical-source-hash: fnv1a64:11c0a9dba3a27832
% canonical-lessons: SA-C415-srotah, SA-C415-avartah, SA-C415-santarati, SA-C415-param, SA-C415-setubandhah

\chapter{River Currents}
\label{ch:sa-river-currents}
\hlchaptermodality{415}

\begin{chapteropening}
\textbf{By the end of this chapter:} \emph{I can say how we got across a river.}

Complete the last lesson of chapter 415: I can say how we got across a river.
\end{chapteropening}

\section[\texorpdfstring{\sk{स्रोतः} (srotaḥ)}{स्रोतः (srotaḥ)}]{\sk{स्रोतः} (srotaḥ) — a flowing current; a stream}

\label{lesson:SA-C415-srotah}



\begin{quote}
Before the new one: say the Sanskrit for unfamiliar, then the Sanskrit for a newcomer.
\end{quote}

\subsection*{You'll want to know: \sk{स्रोतः}}
\textbf{\sk{स्रोतः}} — \emph{srotaḥ} — \textquotedblleft{}a flowing current; a stream\textquotedblright{}.

\begin{quote}
\sk{वर्षायां} \sk{नद्याः} \sk{स्रोतः} \sk{वर्धते।}
\end{quote}

\begin{quote}
\emph{varṣāyāṃ nadyāḥ srotaḥ vardhate.}
\end{quote}

\begin{quote}
In the rains the river's current grows.
\end{quote}

A neuter noun.

\begin{cousinweb}
\textbf{\sk{स्रोतः}} (\emph{srotaḥ}) comes from the root √sru, \textquotedblleft{}to flow\textquotedblright{}; MW gives \emph{srotas} as the current of a river, a stream (MW).
\end{cousinweb}

\begin{grammarlens}[title={What you've built}]
One more word for getting around.
\end{grammarlens}

\subsection*{Your turn}
\begin{itemize}
\raggedright
  \item \emph{Say it:} \emph{srotaḥ}
  \item \emph{Say it:} \emph{srotaḥ}, once more
  \item \emph{Say it:} \emph{varṣāyāṃ nadyāḥ srotaḥ vardhate.} — In the rains the river's current grows.
  \item \emph{Recall:} say the Sanskrit for unfamiliar, then the Sanskrit for a newcomer, then say \emph{srotaḥ} again
\end{itemize}

\begin{tcolorbox}[breakable,colback=teal!4,colframe=teal!35!black,title={Before you move on}]
What does \sk{स्रोतः} mean? (\textquotedblleft{}A flowing current; a stream\textquotedblright{}.) Say it once more.
\end{tcolorbox}
\section[\texorpdfstring{\sk{आवर्तः} (āvartaḥ)}{आवर्तः (āvartaḥ)}]{\sk{आवर्तः} (āvartaḥ) — a whirlpool}

\label{lesson:SA-C415-avartah}



\begin{quote}
Before the new one: say the Sanskrit for a newcomer, then the Sanskrit for a flowing current.
\end{quote}

\subsection*{You'll want to know: \sk{आवर्तः}}
\textbf{\sk{आवर्तः}} — \emph{āvartaḥ} — \textquotedblleft{}a whirlpool\textquotedblright{}.

\begin{quote}
\sk{नद्याः} \sk{मध्ये} \sk{आवर्तः} \sk{अस्ति।}
\end{quote}

\begin{quote}
\emph{nadyāḥ madhye āvartaḥ asti.}
\end{quote}

\begin{quote}
There is a whirlpool in the middle of the river.
\end{quote}

A masculine noun.

\begin{cousinweb}
\textbf{\sk{आवर्तः}} (\emph{āvartaḥ}) is built from \emph{ā} and the root √vṛt, \textquotedblleft{}to turn\textquotedblright{}: a turning round; MW gives \emph{āvarta} as a whirlpool (MW).
\end{cousinweb}

\begin{grammarlens}[title={What you've built}]
One more word for getting around.
\end{grammarlens}

\subsection*{Your turn}
\begin{itemize}
\raggedright
  \item \emph{Say it:} \emph{āvartaḥ}
  \item \emph{Say it:} \emph{āvartaḥ}, once more
  \item \emph{Say it:} \emph{nadyāḥ madhye āvartaḥ asti.} — There is a whirlpool in the middle of the river.
  \item \emph{Recall:} say the Sanskrit for a newcomer, then the Sanskrit for a flowing current, then say \emph{āvartaḥ} again
\end{itemize}

\begin{tcolorbox}[breakable,colback=teal!4,colframe=teal!35!black,title={Before you move on}]
What does \sk{आवर्तः} mean? (\textquotedblleft{}A whirlpool\textquotedblright{}.) Say it once more.
\end{tcolorbox}
\section[\texorpdfstring{\sk{सन्तरति} (santarati)}{सन्तरति (santarati)}]{\sk{सन्तरति} (santarati) — swims across; crosses over}

\label{lesson:SA-C415-santarati}



\begin{quote}
Before the new one: say the Sanskrit for a flowing current, then the Sanskrit for a whirlpool.
\end{quote}

\subsection*{You'll want to know: \sk{सन्तरति}}
\textbf{\sk{सन्तरति}} — \emph{santarati} — \textquotedblleft{}swims across; crosses over\textquotedblright{}.

\begin{quote}
\sk{बालः} \sk{नदीं} \sk{सन्तरति।}
\end{quote}

\begin{quote}
\emph{bālaḥ nadīṃ santarati.}
\end{quote}

\begin{quote}
The boy swims across the river.
\end{quote}

Past (he …): \sk{सः} \sk{नदीं} \sk{सन्तीर्णवान्}, \emph{saḥ nadīṃ santīrṇavān}.

\begin{cousinweb}
\textbf{\sk{सन्तरति}} (\emph{santarati}) is built from \emph{sam} and the root √tṝ, \textquotedblleft{}to cross\textquotedblright{}, which also gives \emph{tarati}, \textquotedblleft{}swims, crosses\textquotedblright{}; MW gives \emph{saṃ-tṝ} as to cross or pass through, with the participle \emph{saṃtīrṇa} (MW).
\end{cousinweb}

\begin{grammarlens}[title={What you've built}]
One more word for getting around.
\end{grammarlens}

\subsection*{Your turn}
\begin{itemize}
\raggedright
  \item \emph{Say it:} \emph{santarati}
  \item \emph{Say it:} \emph{santarati}, once more
  \item \emph{Say it:} \emph{bālaḥ nadīṃ santarati.} — The boy swims across the river.
  \item \emph{Recall:} say the Sanskrit for a flowing current, then the Sanskrit for a whirlpool, then say \emph{santarati} again
\end{itemize}

\begin{tcolorbox}[breakable,colback=teal!4,colframe=teal!35!black,title={Before you move on}]
What does \sk{सन्तरति} mean? (\textquotedblleft{}Swims across; crosses over\textquotedblright{}.) Say it once more.
\end{tcolorbox}
\section[\texorpdfstring{\sk{पारम्} (pāram)}{पारम् (pāram)}]{\sk{पारम्} (pāram) — the far bank}

\label{lesson:SA-C415-param}



\begin{quote}
Before the new one: say the Sanskrit for a whirlpool, then the Sanskrit for swims across.
\end{quote}

\subsection*{You'll want to know: \sk{पारम्}}
\textbf{\sk{पारम्}} — \emph{pāram} — \textquotedblleft{}the far bank\textquotedblright{}.

\begin{quote}
\sk{नौका} \sk{नद्याः} \sk{पारं} \sk{गच्छति।}
\end{quote}

\begin{quote}
\emph{naukā nadyāḥ pāraṃ gacchati.}
\end{quote}

\begin{quote}
The boat goes to the far bank of the river.
\end{quote}

A neuter noun.

\begin{cousinweb}
\textbf{\sk{पारम्}} (\emph{pāram}) comes from the root √pṛ, \textquotedblleft{}to bring across\textquotedblright{}; MW gives \emph{pāra} as the further bank or shore, the opposite side; it lives on as Hindi \emph{pār}, \textquotedblleft{}the further bank\textquotedblright{} (MW; CDIAL 8100).
\end{cousinweb}

\begin{grammarlens}[title={What you've built}]
One more word for getting around.
\end{grammarlens}

\subsection*{Your turn}
\begin{itemize}
\raggedright
  \item \emph{Say it:} \emph{pāram}
  \item \emph{Say it:} \emph{pāram}, once more
  \item \emph{Say it:} \emph{naukā nadyāḥ pāraṃ gacchati.} — The boat goes to the far bank of the river.
  \item \emph{Recall:} say the Sanskrit for a whirlpool, then the Sanskrit for swims across, then say \emph{pāram} again
\end{itemize}

\begin{tcolorbox}[breakable,colback=teal!4,colframe=teal!35!black,title={Before you move on}]
What does \sk{पारम्} mean? (\textquotedblleft{}The far bank\textquotedblright{}.) Say it once more.
\end{tcolorbox}
\section[\texorpdfstring{\sk{सेतुबन्धः} (setubandhaḥ)}{सेतुबन्धः (setubandhaḥ)}]{\sk{सेतुबन्धः} (setubandhaḥ) — a causeway; a dam}

\label{lesson:SA-C415-setubandhah}



\begin{quote}
Before the new one: say the Sanskrit for swims across, then the Sanskrit for the far bank.
\end{quote}

\subsection*{You'll want to know: \sk{सेतुबन्धः}}
\textbf{\sk{सेतुबन्धः}} — \emph{setubandhaḥ} — \textquotedblleft{}a causeway; a dam\textquotedblright{}.

\begin{quote}
\sk{ग्रामस्य} \sk{समीपे} \sk{नद्यां} \sk{सेतुबन्धः} \sk{अस्ति।}
\end{quote}

\begin{quote}
\emph{grāmasya samīpe nadyāṃ setubandhaḥ asti.}
\end{quote}

\begin{quote}
Near the village there is a causeway across the river.
\end{quote}

A masculine noun.

\begin{cousinweb}
\textbf{\sk{सेतुबन्धः}} (\emph{setubandhaḥ}) is built from \emph{setu}, \textquotedblleft{}a bridge, a dam\textquotedblright{}, and \emph{bandha}, \textquotedblleft{}a binding\textquotedblright{}, from the root √bandh, \textquotedblleft{}to bind\textquotedblright{}; MW gives the compound as the making of a causeway, a dam or a bridge (MW).
\end{cousinweb}

\begin{grammarlens}[title={What you've built}]
One more word for getting around.
\end{grammarlens}

\subsection*{Your turn}
\begin{itemize}
\raggedright
  \item \emph{Say it:} \emph{setubandhaḥ}
  \item \emph{Say it:} \emph{setubandhaḥ}, once more
  \item \emph{Say it:} \emph{grāmasya samīpe nadyāṃ setubandhaḥ asti.} — Near the village there is a causeway across the river.
  \item \emph{Recall:} say the Sanskrit for swims across, then the Sanskrit for the far bank, then say \emph{setubandhaḥ} again
\end{itemize}

\begin{tcolorbox}[breakable,colback=teal!4,colframe=teal!35!black,title={Before you move on}]
What does \sk{सेतुबन्धः} mean? (\textquotedblleft{}A causeway; a dam\textquotedblright{}.) Say it once more.
\end{tcolorbox}
